% =====================================================================
\section{Distance certification and quantum parameter analysis}
\label{sec:cert}
% =====================================================================

\subsection{The coset-distance theorem}

The distance formulas of Theorems~\ref{thm:R2} and~\ref{thm:R3} look like
minima over a set difference. The following theorem says that they are
something better: the minimum weight of an ordinary linear code.

\begin{theorem}[coset-distance characterisation]\label{thm:cosetdistance}
Let $D\subseteq\F_q^{\Nlen}$ be $\F_q$-linear, $H_E=D\cap D^{\perp_E}$ and
$K=\Nlen-k-h_E$. Then, for routes R2 and R3,
\[
  d=\min\bigl\{\wt(x):x\in D^{\perp_E}\setminus\Hull_E(D)\bigr\}
\]
is the \emph{minimum weight of the quotient code}
\[
  D^{\perp_E}/\Hull_E(D),
\]
a linear code of length $\Nlen$ and dimension
$\dim D^{\perp_E}-\dim\Hull_E(D)=(\Nlen-k)-h_E=K$. Equivalently, $d$ is the
minimum weight of the set of nonzero cosets of $\Hull_E(D)$ in $D^{\perp_E}$.
In particular, for any integer $\delta\ge1$,
\begin{equation}\label{eq:criterion}
  d\ge\delta
  \iff
  \text{every }x\in D^{\perp_E}\text{ with }\wt(x)<\delta
  \text{ lies in }\Hull_E(D).
\end{equation}
\end{theorem}

\begin{proof}
By Theorem~\ref{thm:parameters}, for R2 and R3
$d$ is the minimum of $\swt$ over $G^{\omegaperp}\setminus G$ with $G=V$.
Since $\swt(a|b)\ge\max(\wt a,\wt b)$ and since
$(a|0)\in V^{\omegaperp}\setminus V$ whenever $a\in D^{\perp_E}\setminus D$,
that symplectic minimum equals
$\min\{\wt(x):x\in D^{\perp_E}\setminus D\}$; and for $x\in D^{\perp_E}$ one
has $x\in D$ iff $x\in D\cap D^{\perp_E}=H_E$, so
$D^{\perp_E}\setminus D=D^{\perp_E}\setminus H_E$ as sets.

Now, the quotient group $D^{\perp_E}/H_E$ is an $\F_q$-vector space of
dimension $\dim D^{\perp_E}-\dim H_E=(\Nlen-k)-h_E=K$, and its nonzero elements
are exactly the cosets $x+H_E$ with $x\in D^{\perp_E}\setminus H_E$. The weight
of a coset is well defined as
$\wt(x+H_E)=\min\{\wt(y):y\in x+H_E\}$, and its minimum over nonzero cosets is
exactly $\min\{\wt(x):x\in D^{\perp_E}\setminus H_E\}$: indeed
$\wt(x+H_E)\le\wt(x)$ gives one inequality, and conversely any $y\in x+H_E$
with $y\ne0$ may be substituted for $x$. This is the assertion. The criterion
\eqref{eq:criterion} is the contrapositive restatement.
\end{proof}

\begin{remark}
For route R1 only the upper bound of Theorem~\ref{thm:R1cap} is available, and
Theorem~\ref{thm:cosetdistance} is \emph{not} claimed for R1: the group
$\mathcal C_T$ is semilinear, the set difference is not a quotient by an
$\F_q$-subspace of the ambient space in the same way, and we do not use it.
\end{remark}

\subsection{The coset Singleton bound}

\begin{theorem}[coset Singleton bound]\label{thm:cosetsingleton}
Let $D\subseteq\F_q^{\Nlen}$ be $\F_q$-linear and let $K=\Nlen-k-h_E$ be the
dimension of the quantum code produced by route R2 or R3. Then
\begin{equation}\label{eq:cosetsingleton}
  K+d\le \Nlen+1 .
\end{equation}
\end{theorem}

\begin{proof}
Consider the quotient code $\mathcal Q=D^{\perp_E}/\Hull_E(D)$ of
Theorem~\ref{thm:cosetdistance}: it has length $\Nlen$, dimension $K$ and
minimum weight $d$. Let $S\subseteq\{1,\dots,\Nlen\}$ be any set of
$\Nlen-d+1$ coordinates and let $\pi_S$ be the projection onto those
coordinates. If $\pi_S$ were not injective on $\mathcal Q$, there would be a
nonzero coset $x+H_E$ with $\pi_S(x+H_E)=0$, i.e. an element
$y\in x+H_E$ supported outside $S$, i.e. supported on the remaining
$d-1$ coordinates. Then $\wt(y)\le d-1<d$, contradicting the definition of $d$
as the minimum weight of a nonzero element of $\mathcal Q$. Hence $\pi_S$ is
injective, so $q^K=|\mathcal Q|\le q^{\Nlen-d+1}$ and therefore
$K\le\Nlen-d+1$.
\end{proof}

\begin{proposition}[the coset Singleton bound is new]\label{prop:newsbound}
The bound \eqref{eq:cosetsingleton} is not the entanglement-assisted Singleton
bound, and neither implies the other.
\end{proposition}

\begin{proof}
The entanglement-assisted Singleton bound reads
$\Nlen-K+c\ge 2(d-1)$, i.e. $K\le\Nlen+c-2d+2$. For the code
$[[20,2,10;18]]_4$ constructed in Section~\ref{sec:cmp} one has $\Nlen=20$,
$K=2$, $d=10$, $c=18$: the coset Singleton bound gives $2\le 20-10+1=11$, which
holds, while the EA Singleton bound gives
$2\le 20+18-20+2=20$, which is much weaker. Conversely, for a code with small
$c$ and small $d$, say $[[24,20,2;0]]_4$ ($\Nlen=24$, $K=20$, $d=2$, $c=0$), the
coset Singleton bound gives $20\le 24-2+1=23$, while the EA Singleton bound
gives $20\le 24+0-4+2=22$ --- the EA bound is now the stronger one. Hence
neither implies the other.
\end{proof}

\subsection{The feasibility inequality}

\begin{theorem}[feasibility inequality]\label{thm:feasibility}
Suppose a code produced by route R3 from $D\subseteq\F_q^{\Nlen}$ has
$d\ge4$, $c\le c_0$ and $K\ge K_0$. Then
\begin{equation}\label{eq:feasibility}
  2k\le \Nlen+c_0-K_0 .
\end{equation}
\end{theorem}

\begin{proof}
By Theorem~\ref{thm:R3}, $c=k-h_E$ and $K=\Nlen-k-h_E$. The hypotheses
$c\le c_0$ and $K\ge K_0$ read $h_E\ge k-c_0$ and
$h_E\le \Nlen-k-K_0$. Hence $k-c_0\le \Nlen-k-K_0$, i.e.
$2k\le\Nlen+c_0-K_0$.
\end{proof}

\begin{remark}[what the inequality does and does not say]
The hypothesis $d\ge4$ enters only through the \emph{motivation} for imposing
the constraints, not through the algebra: the implication
``$c\le c_0$ and $K\ge K_0$ $\Rightarrow$ \eqref{eq:feasibility}'' holds for
every code, whatever its distance. Inequality \eqref{eq:feasibility} is one-way:
it is a necessary condition, not a sufficient one, and no converse is claimed.

Its content is the trade-off behind the entanglement budget: cheap entanglement
means a large Euclidean hull, and a large Euclidean hull means few logical
qudits. Combined with the criterion \eqref{eq:criterion}, it converts the
question ``is there a family with $d\ge4$, bounded $c$ and $K\ge2$?'' into two
explicit conditions on $D$: (i) $k$ must be close to $\Nlen/2$, and (ii) every
vector of weight at most $3$ in $D^{\perp_E}$ must lie in $\Hull_E(D)$.
\end{remark}

\subsection{Bookkeeping identities, and an erratum}

\begin{theorem}[bookkeeping]\label{thm:bookkeeping}
For every code produced by any of the three routes from $C\le\Rn$ with
$D=\varphi_n(C)$:
\begin{enumerate}[label=\textnormal{(\alph*)},leftmargin=*]
\item $k=\dim_{\F_q}D$ is even, and $\Nlen=4tn$ is divisible by $4$;
\item $K-c=\Nlen-2k\equiv0\pmod4$; in particular $K\equiv c\pmod4$;
\item $c+K=\Nlen-2h_E$;
\item $c\ge 2k-\Nlen$.
\end{enumerate}
\end{theorem}

\begin{proof}
(a) Under the Morita correspondence \eqref{eq:morita}, $D$ is the image of
$I_W$ with $W\le\Bq^2$, and $\dim_{\F_q}I_W=2\dim_{\F_q}W$; hence $k$ is even.
(b) $K-c=(\Nlen-k-h_E)-(k-h_E)=\Nlen-2k$, which is divisible by $4$ by (a)
and $\Nlen=4tn$. (c) is immediate from the two defining identities. (d) From
$c=k-h_E$ and $h_E\le\min(k,\Nlen-k)$ (as $h_E$ is the dimension of a subspace
of both $D$ and $D^{\perp_E}$) we get $c\ge k-(\Nlen-k)=2k-\Nlen$.
\end{proof}

\begin{remark}[erratum]
An earlier version of this project stated the identity as $c+K=\Nlen-k$. That
is wrong: it would force $h_E=k/2$, which is false in general, and it is
contradicted by $23$ of the $25$ rows of the earlier paper's own parameter
table. The correct identity is Theorem~\ref{thm:bookkeeping}(c). We record this
as an erratum rather than repair it silently; no \emph{parameter} computed in
the earlier work is affected, because $K$ and $c$ were always computed from
$k$ and $h_E$ directly.
\end{remark}

\subsection{Certification methodology}

Every distance reported in this paper is either \textbf{exact} or a
\textbf{proved lower bound}, and the two are never conflated. Three
independent methods are employed across the certification; their applicability
is recorded per code, and no single code is certified by all three. Each
headline code has its exact distance certified by exactly one of the two exact
methods: M1 where the enumeration completes, and otherwise M2, which is run
precisely when M1 cannot finish. Every headline code is additionally
cross-checked by the ambient quantum search (M3), which for most codes yields
a proved lower bound rather than an exact value (see
Appendix~\ref{app:tables} and the supplementary certification data). No code
is certified by all three.

\begin{description}
\item[M1 --- complete enumeration.] Enumerate the relevant classical difference
  set (the coset code of Theorem~\ref{thm:cosetdistance}, or the classical
  difference set of R2) completely and take the minimum weight. This is exact
  and complete whenever $q^{\dim}\le3\cdot10^6$.
\item[M2 --- increasing-weight support search.] A different algorithm: for each
  support of size $w=1,2,3,\dots$ in increasing order, solve the left-kernel
  problem asking whether the difference set contains a vector supported inside
  it. The first $w$ at which a solution exists is the distance. Because the
  algorithm differs structurally from M1, agreement of the two is a genuine
  cross-check.
\item[M3 --- ambient quantum search.] Search the ambient space
  $\F_q^{2\Nlen}$ in increasing symplectic weight for a vector $v$ with
  $v\in V^{\omegaperp}$ and $v\notin V$. This certifies the distance directly in
  the quantum language, without passing through
  Theorem~\ref{thm:cosetdistance}. It is complete only up to the weight reached,
  and is reported as a lower bound otherwise.
\end{description}

In addition, for every certified code: the stabilizer group is verified abelian
on an $\F_p$-basis (Theorem~\ref{thm:fpbasis}); $K$ and $c$ are recomputed from
$k$ and $h_E$; and --- crucially --- \emph{the certified code is the exact code
found by the search}, the image basis being carried through the whole pipeline.
An earlier version of this project certified a re-selected code with the same
invariants instead of the reported one; that defect was found and fixed, and the
current pipeline cannot reproduce it.

\subsection{Parameter analysis}

The computation covers $q\in\{4,8,9\}$ with $tn\le8$ (42 families), of which
$23$ are enumerated exhaustively, $4$ are counted but not enumerated, and $15$
are blocked by the component enumeration. Ten families that earlier work could
only sample are now exhaustive. In total $120\,721$ codes have computed
invariants and $5\,055$ have a computed distance, of which $4\,418$ are exact
and $637$ are proved lower bounds. Every one of the $5\,055$ satisfies the
entanglement-assisted Singleton and Hamming bounds, and the coset Singleton
bound \eqref{eq:cosetsingleton} (0 violations in each case).

\begin{proposition}[frontier, computational]\label{prop:frontier}
Among the $5\,055$ codes with a distance there are $368$ distinct tuples
$(\Nlen,K,d,c)$; $88$ of them are non-dominated and $280$ are dominated, where
$(N,K,d,c)$ dominates $(N',K',d',c')$ iff
$N\le N'$, $K\ge K'$, $d\ge d'$, $c\le c'$ with at least one strict
inequality, the comparison being made inside one alphabet.
\end{proposition}

\begin{remark}[what the frontier is and is not]
Proposition~\ref{prop:frontier} is a statement about the codes \emph{scanned},
not about all codes obtainable from the construction: the families that were
counted but not enumerated, and the fifteen that were blocked, are excluded, and
they are listed in the supplementary material. No tuple outside the scanned
families is claimed to be dominated.
\end{remark}

Table~\ref{tab:headline} lists the headline codes, each certified by an
applicable exact method and cross-checked by the ambient quantum search. Table~\ref{tab:cert} in Appendix~A lists the full certification table.

\begin{table}[t]
\centering
\caption{Headline codes. $S$-def is the entanglement-assisted Singleton defect
$\Nlen-K+c-2(d-1)$. All distances exact except the last row of the upper
block, which is a proved lower bound.}
\label{tab:headline}
\small
\begin{tabular}{@{}lllrrrrrr@{}}
\toprule
code & family $(q,t,n)$ & route & $k$ & $h_E$ & $K$ & $d$ & $c$ & $S$-def\\
\midrule
$[[24,8,4]]_9$     & $(9,1,6)$ & R2-SO & 8  & 8 & 8  & 4  & 0  & 10\\
$[[20,4,4]]_4$     & $(4,1,5)$ & R2-SO & 8  & 8 & 4  & 4  & 0  & 10\\
$[[20,4,4]]_4$     & $(4,1,5)$ & R2-DC & 12 & 8 & 4  & 4  & 0  & 10\\
$[[16,4,3]]_4$     & $(4,1,4)$ & R2-SO & 6  & 6 & 4  & 3  & 0  & 8\\
$[[16,4,3]]_4$     & $(4,1,4)$ & R2-DC & 10 & 6  & 4 & 3  & 0  & 8\\
$[[20,2,10;18]]_4$ & $(4,1,5)$ & R3    & 18 & 0 & 2  & 10 & 18 & 18\\
$[[24,2,10;18]]_9$ & $(9,1,6)$ & R3    & 20 & 2 & 2  & 10 & 18 & 22\\
$[[20,4,9;16]]_9$  & $(9,1,5)$ & R3    & 16 & 0 & 4  & 9  & 16 & 16\\
$[[24,4,9;20]]_4$  & $(4,1,6)$ & R3    & 20 & 0 & 4  & 9  & 20 & 24\\
$[[16,2,8;14]]_9$  & $(9,1,4)$ & R3    & 14 & 0 & 2  & 8  & 14 & 14\\
$[[16,2,7;10]]_9$  & $(9,1,4)$ & R3    & 12 & 2 & 2  & 7  & 10 & 12\\
$[[12,2,6;10]]_4$  & $(4,1,3)$ & R3    & 10 & 0 & 2  & 6  & 10 & 10\\
$[[24,4,6;20]]_4$  & $(4,2,3)$ & R3    & 20 & 0 & 4  & 6  & 20 & 30\\
$[[16,4,5;8]]_4$   & $(4,1,4)$ & R3    & 10 & 2 & 4  & 5  & 8  & 12\\
$[[24,10,4;2]]_9$  & $(9,1,6)$ & R3    & 8  & 6 & 10 & $\ge4$ & 2 & 10\\
\midrule
$[[24,6,2;2]]_9$   & $(9,2,3)$ & R3 & 10 & 8  & 6 & 2 & 2 & 18\\
$[[24,14,2]]_9$    & $(9,2,3)$ & R1 & 10 & --- & 14 & 2 & 0 & 8\\
$[[32,8,2;4]]_4$   & $(4,2,4)$ & R3 & 14 & 10 & 8 & 2 & 4 & 26\\
\bottomrule
\end{tabular}
\end{table}

The last three rows are the non-principal-hull witnesses of
Theorem~\ref{thm:nonprincipal}: their hulls have $\MU_{\Bq}=3$. They are
included for their structural significance, not for their parameters
($d=2$). The middle one of the three is a \emph{stabilizer} code produced by
route~R1, which is the only route available there because
Theorem~\ref{thm:obstruction} blocks the Euclidean one at $(q,t)=(9,2)$.

\subsection{Bounds satisfied}

\begin{proposition}[bounds]\label{prop:bounds}
All $5\,055$ codes with a computed distance satisfy
\begin{enumerate}[label=\textnormal{(\alph*)},leftmargin=*]
\item the entanglement-assisted Singleton bound $\Nlen-K+c\ge2(d-1)$;
\item the entanglement-assisted Hamming bound with effective dimension $K-c$:
  $\bigl(\sum_{i=0}^{t}\binom{\Nlen}{i}(q^2-1)^i\bigr)q^{K-c}\le q^{\Nlen}$
  with $t=\lfloor(d-1)/2\rfloor$;
\item the coset Singleton bound $K+d\le\Nlen+1$.
\end{enumerate}
A lower-bound distance is tested with its proved value.
\end{proposition}

\subsection{The entanglement budget}

Table~\ref{tab:budget} is the measured answer to the question behind the
criterion of Theorem~\ref{thm:feasibility}: what is the least entanglement at
which a given distance becomes available.

\begin{table}[t]
\centering
\caption{Minimum entanglement $c$ for each distance, over the whole search,
with the attaining code. The distance range is $d=2,\dots,8$: these are the
distances for which Part~4(c) of the Phase-7 log records a minimum-entanglement
attaining code. The two larger distances of the search, $d=9$ and $d=10$, are
attained only at $c=16$ and $c=18$ (Table~\ref{tab:headline}).}
\label{tab:budget}
\begin{tabular}{@{}clll@{}}
\toprule
$d$ & minimal $c$ & attaining code & status of $d$\\
\midrule
2 & 0  & $[[24,20,2;0]]_4$  & exact\\
3 & 0  & $[[24,12,3;0]]_9$  & exact\\
4 & 0  & $[[24,8,4;0]]_9$   & exact\\
5 & 2  & $[[24,6,5;2]]_9$   & proved lower bound\\
6 & 10 & $[[12,2,6;10]]_4$  & exact\\
7 & 10 & $[[16,2,7;10]]_9$  & exact\\
8 & 14 & $[[16,2,8;14]]_9$  & exact\\
\bottomrule
\end{tabular}
\end{table}

Reading Table~\ref{tab:budget} with Theorem~\ref{thm:cosetdistance}: $d=4$ is
available with \emph{no} entanglement at all ($[[24,8,4]]_9$, a stabilizer
code), while $d\ge5$ costs entanglement, and the cost grows steeply. With the
entanglement bounded by $c\le4$ the best attainable distance is $d=4$, attained
with $K=10$ and $c=2$ by $[[24,10,4;2]]_9$ --- a proved lower bound, and the
frontier point for the criterion.

\subsection{Worked example VI: a full certification}

% The optional argument is braced: LaTeX's [..] scanner stops at the first
% unbraced ']', and the quantum-code tuple contains ']]'.
\begin{example}[{certifying $[[20,2,10;18]]_4$}]\label{ex:certify}
Take the family $(q,t,n)=(4,1,5)$, so $\Nlen=4\cdot1\cdot5=20$. The search
finds a cyclic code $C\le\Rn$ whose Gray image $D=\varphi_5(C)\subseteq
\F_4^{20}$ has
\[
  k=\dim_{\F_4}D=18,\qquad h_E=\dim\bigl(D\cap D^{\perp_E}\bigr)=0 .
\]
Route R3 (Theorem~\ref{thm:R3}) then gives
\[
  \begin{gathered}
    c=k-h_E=18,\qquad K=\Nlen-k-h_E=20-18-0=2,\\
    [[20,2,d;18]]_4 .
  \end{gathered}
\]
The distance. By Theorem~\ref{thm:cosetdistance}, $d$ is the minimum weight of
the quotient code $\mathcal Q=D^{\perp_E}/\Hull_E(D)$; here $h_E=0$, so
$\Hull_E(D)=\{0\}$ and $\mathcal Q=D^{\perp_E}$, a linear $[20,2]_4$ code. Its
$4^2=16$ codewords are enumerated completely (method M1), and the minimum
weight found is $10$. The support search (M2) is run only when M1 cannot
finish, so it was not run here and contributes nothing; the ambient quantum
search (M3) independently certifies $d\ge3$ as a lower bound, consistent with
this value. Hence $d=10$ exactly and
\[
  [[20,2,10;18]]_4 .
\]
The three bound checks:
\begin{enumerate}[label=\textnormal{(\roman*)},leftmargin=*]
\item entanglement-assisted Singleton:
  $\Nlen-K+c=20-2+18=36\ge2(d-1)=18$, with defect $18$;
\item entanglement-assisted Hamming with effective dimension $K-c=-16$:
  $\bigl(\sum_{i=0}^{4}\binom{20}{i}(4^2-1)^i\bigr)4^{-16}\le4^{20}$;
\item coset Singleton (Theorem~\ref{thm:cosetsingleton}):
  $K+d=2+10=12\le\Nlen+1=21$.
\end{enumerate}
For comparison, the entanglement-assisted Singleton bound alone only gives
$K\le\Nlen+c-2d+2=20$, so the coset Singleton bound is the binding one here
(Proposition~\ref{prop:newsbound}).

Finally, the bookkeeping identities of Theorem~\ref{thm:bookkeeping}:
$k=18$ is even; $K-c=2-18=-16=\Nlen-2k=20-36$ is divisible by $4$; and
$c+K=20=\Nlen-2h_E$, in accordance with the corrected identity --- the
earlier formula $c+K=\Nlen-k$ would give $20=2$, which is false.

No code found anywhere in the search has $d>10$; together with the
$[[24,2,10;18]]_9$ of Example~\ref{ex:routes}, this $[[20,2,10;18]]_4$ is
one of the two largest-distance codes of the search, and it is also the
most entanglement-hungry ($c=18$ out of $\Nlen=20$). That trade-off is
the content of Theorem~\ref{thm:feasibility}: $d\ge4$ together with
$K\ge2$ forces $2k\le\Nlen+c_0-K_0$, i.e. large $d$ and moderate $K$ cannot
both be had at small $c$.
\end{example}
